\section{Finite pressure expansion and power count}\label{sec:expansion}

We now estimate the finite pressure polynomial that represents the
increment in Proposition~\ref{prop:increment}. The moment bounds control
its high-degree terms. For degrees two, three, and four, we must first
sum the signed scalar coefficients: the required saving appears only
after this sum and, on a size path, after applying the size derivative.

All estimates are on the stopped set. On direct strength paths we
restrict to \(Y\le x^\kappa\), since
Lemma~\ref{lem:statistics-coarse-switch} treats its complement. On size
paths we use the cap of Corollary~\ref{cor:statistics-cap}. Group masses
up to a fixed multiple of \(Y\) into one lower band, assigned scale
\(r=Y\), and use dyadic bands \(v\asymp r\) above it. Fixed changes in
these endpoints affect only constants. Write
\[
\delta=Q^*-S_v,\qquad f=S_v-q(v),\qquad \Delta=\delta+f.
\]
Here \(S_v\) is the untilted mean, used as the same deterministic
centering in both ensembles. When specified, \(Q^*\) is replaced by the
masked inner product of posterior means at a common later cut. The
allocation estimate below permits these projected factors. On size
paths we apply it only to factors in the low region of \eqref{eq:D2}.
The bounds on \(Q-q\) and \(S-q\) give the same envelope bound for
\(\delta\) by the triangle inequality; on the lower band it is \(O_p(Y)\).

\subsection{The pressure polynomial and its derivatives}

Choose a fixed Taylor degree \(J_0\), allowing different degrees for
size and direct variations. The size degree can be chosen from
\eqref{eq:D11} before any final decrease of \(p_*\). Take the size
stencil order at least this degree. Lemma~\ref{found:size-increment}
expresses the change from \(m\) to \(m+j\) as a polynomial in \(j\).
The stencil extracts its linear coefficient, which is the per-root
expectation of the whole symbol \eqref{eq:A6}. The path derivative
adds \(\partial_\theta\) of that symbol. Thus the operator for the size
increment is \({\cal O}=1+\partial_\theta\), with the unit acting on
every term. For a direct strength increment the operator is instead
\({\cal O}=\partial_{\log\lambda}\).

These operators act on formal polynomials: on a size path assign
\(\partial_\theta\Delta=-\Delta/6\), and on a direct path hold
\(\Delta\) fixed. This is the rule in Lemma~\ref{found:row-derivative},
not a derivative of a random overlap or of its law. Scalar coefficients
receive their ordinary parameter derivatives. After the exact
cancellation of the singleton cumulant with the linear counterterm,
\eqref{eq:A6} becomes, up to deterministic terms,
\[
 -\tfrac14\sum b\Delta^2+
 \sum_{j=2}^{J_0}\frac1{2^j j!}
   \sum_{\text{pairs}}\kappa_K(\tau_{a_1}\tau_{b_1},\ldots,\tau_{a_j}\tau_{b_j})
                 \prod_{i=1}^j (b\Delta+d)_{a_i b_i}.                \tag{E1}\label{eq:E3}
\]
where the pairs are ordered and off diagonal, and \(\kappa_K\) is their
joint scalar cumulant at fixed topology. We must bound the difference
of per-root positive-law expectations after applying \({\cal O}\) to
\eqref{eq:E3}. The degree counts below always refer to the base pair factors in this
polynomial. Coefficient differentiation may allocate an additional scalar
pair, whose clock weight is treated separately.

For high masses on size paths, \eqref{eq:D15} gives the needed gain.
For low masses, the remaining task is to obtain a coefficient of size
\(\lambda w^{5-j}\) for a product of \(j\) factors at scales at most
\(w\). This is why the next lemma uses that particular power of
\(w\). Its extra ratio for two distinct quadratic vertices records
the smaller split scale. We will obtain these coefficient bounds by
summing the signed scalar contributions of the polynomial; before doing
so, we prove the estimate that converts them into an increment bound.

\subsection{A centered allocation estimate}

\begin{lemma}[Centered allocation estimate]\label{exp:lem-centered-count}
Fix \(j\ge2\) pair factors. Each factor is either \(\delta\), possibly
projected, or a deterministic function bounded by
\(O_\diamond(e(v))\); use \(e=Y\) on the lower band. Fix the bands, their
order, and every projection cut. Each projection cut lies after the
split of its pair. Let \(s=1,\ldots,l\) index the distinct vertices to
which these factors are assigned, with multiplicities \(k_s\) and band
scales \(r_s\), and put \(w=\max_s r_s\). Suppose that the deterministic
weights are common to both ensembles and, on each of boundedly many
topology shapes, are bounded by
\[
 O_\diamond(\lambda w^{5-j})\ \prod_{s=1}^l dv_s/w,                    \tag{E2}\label{eq:E1}
\]
times measures of bounded volume in the other integration variables.
These are product majorants, so deterministic restrictions may be
retained. If \(j=l=2\), assume in addition the factor
\(O_\diamond(\min_s r_s/w)\). More explicitly, if \(H\) is this
deterministic weight and \(T_i^{\mathrm{tilt}},T_i^0\) are the pair
factors in the two ensembles, we estimate
\[
 \int H\left\{\nu^{\mathrm{tilt}}
       \left[\prod_{i=1}^jT_i^{\mathrm{tilt}}\right]
       -\nu^0\left[\prod_{i=1}^jT_i^0\right]\right\}.
\]
The weight, deterministic factors, cuts, and centering \(S\) are common;
only the positive law and the masked overlap vary. The absolute value
of this difference is bounded by
\[
 O_\diamond\big(\lambda Y^5(x/Y+g_B)\big)+O(x^K)                       \tag{E3}\label{eq:E2}
\]
where \(K\) can be arbitrarily large. Indeed the two contributions are
\[
 \lambda Y^5\frac{x}{Y}=x^5\lambda^{1-2s_*}Z^4,
 \qquad
 \lambda Y^5g_B=x^5\lambda^{(1-\eta)/2-s_*}Z^{7/2}.
\]
Both powers of \(\lambda\) are positive. Since
\(\mathbf E_\gamma Z^4\lesssim_\circ1\), the averaged bound is
\(O(x^{5+c'})\) for some \(c'>0\), after taking the allowed losses smaller
than these gains.
\end{lemma}

\begin{proof}
Set \(a_s=r_s/Y\), \(A=w/Y\), and \(m_{r_s}=Y/a_s\). The ratios
\(a_s\) and \(A\) are at least one. At a vertex carrying \(k\) factors,
\eqref{eq:D2} and the integrated envelope bound give
\[
 P_k=m_{r_s}^{\min(2,k)}r_s^{k-\min(2,k)}
\]
after integration in \(dv_s/r_s\). Thus one or two factors use the
first or second envelope moment, and each further factor costs \(r_s\).
Changing \(dv_s/w\) to \((r_s/w)dv_s/r_s\) in \eqref{eq:E1} shows that
the multiplier relative to \(\lambda Y^5\) is
\[
 F=A^{5-j-l}\prod a_s^{1+k_s-2\min(2,k_s)}
\]
with the additional factor \(\min_s a_s/A\) when \(j=l=2\).
The exponent of an individual \(a_s\) is \(0\) for a single factor,
\(-1\) for a double factor, and \(k_s-3\) for \(k_s\ge3\). Consequently
\[
\begin{array}{c|c}
\text{multiplicities}&F\text{ or an upper bound for }F\\ \hline
l=1&A\\
(1,2)&1/a_{\rm duplicate}\\
(1,1)&\min(a_1,a_2)\\
\text{all other cases with }l\ge2&A^{-1}.
\end{array}
\]
If an assigned vertex lies below \(x\), its integration uses only the
fraction \(x/Y\) of the lower band. In this case at least one \(a_s=1\),
and every entry in the table is at most a constant. Raw bounds in the
two ensembles therefore give the first contribution in \eqref{eq:E2}.

We may now assume that all assigned vertices are at least \(x\).
Fully deterministic products cancel, so there is a stochastic factor.
Suppose first that all stochastic factors occur at one vertex \(s\).
At that vertex, replace the overlaps by the common ancestor norm
squared and then by \(2L_v\) on a path through the fork. The alias
errors are controlled by \eqref{eq:D7}. Center the resulting polynomial
at its untilted expectation and apply \eqref{eq:D10}, treating the
other times and restrictions as deterministic coefficients. After
truncating the other factors at their envelopes, the cost is
\(O_\diamond(z_{r_s}^0r_sZ)\) times \(1\) for \(k_s=1\), or times
\(m_{r_s}r_s^{k_s-2}\) for \(k_s\ge2\). Relative to \(P_{k_s}\), this
gives at least the gain \(g_B/a_s\). On the lower band the same gain
follows from conditional averaging and the linear bound after
\eqref{eq:D9} when there is one stochastic power, and from the
integrated \(D_v^2\) bound when there are two or more. Multiplying by
the appropriate entry for \(F\) gives \(O_\diamond(g_B)\).

If stochastic factors occur at several vertices, select two of them
at scales \(aY,bY\). Their integrated second moments about \(S\), from
\eqref{eq:D7} and \eqref{eq:D9}, are bounded by
\(O_\diamond(g_BY^2/a)\) and \(O_\diamond(g_BY^2/b)\).
First apply unconditional Cauchy--Schwarz to the two selected stochastic
factors at fixed masses, and truncate the remaining factors at their
deterministic envelopes. At a selected vertex of scale \(r=aY\)
and multiplicity \(k\ge2\), Cauchy--Schwarz in mass then gives
\[
 \int e(v)^{k-1}D_v\,\frac{dv}{r}
 \lesssim r^{k-2}
    \left(\int e(v)^2\frac{dv}{r}\right)^{1/2}
    \left(\int D_v^2\frac{dv}{r}\right)^{1/2}
 \lesssim_\diamond r^{k-2}m_r
                 (g_BY^2/a)^{1/2}.
\]
Dividing by the raw budget \(P_k=m_r^2r^{k-2}\) gives the factor
\(O_\diamond(\sqrt{g_Ba})\); for \(k=1\), division of
\((\int D_v^2\,dv/r)^{1/2}\) by \(P_1=m_r\) gives the same
factor. Thus the two selected vertices improve the raw product by
\(O_\diamond(g_B\sqrt{ab})\). Since \(\sqrt{ab}\le A\),
this proves the claim whenever \(F\lesssim A^{-1}\).

There are two exceptional multiplicity patterns. For \((1,2)\), let
\(a_d\) be the duplicate scale and \(a_s\) the singleton scale. If both
factors at the duplicate are stochastic, selecting them uses
\(\int D_v^2\,dv/r\) at that vertex. Relative to its raw budget
\(m_r^2\), this gains \(O_\diamond(g_Ba_d)\), so
\(F g_Ba_d=g_B\). If only one is stochastic, selecting it together
with the singleton gives
\[
 Fg_B\sqrt{a_sa_d}=g_B\sqrt{a_s/a_d},
\]
which suffices when \(a_s\lesssim a_d\). In the remaining case the
later singleton is linear and stochastic. Integrating it as a centered
path integral gives \(g_B/A\), so the product with \(F=1/a_d\) is
again at most \(g_B\). For \((1,1)\), both factors must be stochastic.
The lower-band case has \(a_1=a_2=1\) and follows directly from
Cauchy--Schwarz. Otherwise the later linear factor has scale \(A\)
and gives \(g_B/A\); multiplying by
\(F=\min(a_1,a_2)\) gives at most \(g_B\).

For completeness, the last centered integration has a fixed marginal
law even though its fork moves. There is only one earlier stochastic
pair to retain. At most one of its leaves descends through the later
fork, so choose that leaf for the \(L\)-path when it is present. After
the alias replacement the other child of the moving fork has no
observable and integrates to one. The retained pair always splits at
its fixed earlier time. A prior projection merely evaluates these
same two paths at fixed later cuts. If neither endpoint descends
through the moving fork, choose a continuation whose separation from
the retained tree is at another fixed vertex. In either case the
moving time occurs only in the \(L\)-observable and in deterministic
kernels. Apply \eqref{eq:A2}, or its consequence \eqref{eq:D10},
unconditionally under this marginal and then apply H\"older with the
earlier random factor. This is the construction in
Remark~\ref{rem:statistics-earlier-pair}; it does not require an
inequality conditioned on that factor. Deterministic ties belong to
the kernels, and stochastic ties form one polynomial power. The
truncation and tilted-law transfers follow the conventions after
\eqref{eq:D10}. This proves \eqref{eq:E2}.
\end{proof}


\subsection{Taylor remainders and separation of the masses}

On size paths, \eqref{eq:D11} makes each insertion small in every fixed
moment. A sufficiently high fixed degree therefore bounds the Taylor
remainders using the absolute allocation measure. The differentiated
clock weights in \eqref{eq:A9} have absolute integrated norm
\(O_\circ(1)\). Their placements have bounded densities, including when
the added pair attaches to an existing integration variable; there is
no prescribed fork in this pressure calculation. Thus they cost only
a subpower factor.

The same weighted accounting applies to row contacts and normalization
errors. Since \(m\asymp N=x^{-6}\), their contribution is bounded by
\[
 O(m^{-1})\left(1+
   \|\dot b\|_{L^1}+\|\dot h\|_{L^1}
        +\|\dot K\|_{L^1}+\|\dot q\|_{L^1}\right)
       =O_\circ(x^6).
\]
Here the norms denote the absolute mass integrals with the
bounded-density allocation weights from \eqref{eq:A9}. The stencil
has finitely many fixed coefficients. Moreover,
Lemma~\ref{found:size-increment} gives a total-pressure remainder, so
no further factor \(m\) multiplies these errors. This proves that
\eqref{eq:A7}--\eqref{eq:A9} and the stencil reduce the size increment
to \eqref{eq:E3} with the required saving. On direct paths the
differentiated weights carry the additional factor
\(O_\circ(\lambda)\).

For a direct variation, start with the monotone reference
\(K_S=h+bS\). Apply \eqref{eq:A7} with insertions \(b\delta\), expanding
the pressure gradient of \eqref{eq:A9} through order \(J_0-1\).
By \eqref{eq:P8}, the test multiplying the cavity moment has degree
at most one and weight \(O_\circ(\lambda)\). In the true Taylor
remainder distinguish two cases. If at most one insertion has
\(e(v)\ge x^{\delta'}\), at least \(J_0-2\) insertions have smaller
envelope. Their moment bounds from \eqref{eq:D2} supply an arbitrarily
large fixed power of \(x\) once \(\delta'>0\) is fixed and \(J_0\) is
large enough. Choose the truncation losses first so that they do not
use up this gain.

If two insertions have \(e(v)\ge x^{\delta'}\), apply Cauchy--Schwarz
and \eqref{eq:D8} to those insertions, and bound the others uniformly.
The selected vertices have bounded absolute density even if the
insertions are tied. The one-row interpolated laws are comparable to
the genuine laws by \eqref{eq:A7}; the transfer to tilted active
components uses \eqref{eq:A4} and the stopping conventions after
Lemma~\ref{lem:statistics-centered}. Before averaging over \(\gamma\),
the cost is
\[
 O_\diamond\!\left(\lambda x^{-C_0\delta'}z_1^0Y^2Z^2\right)
 =O_\diamond\!\left(
 x^{5-C_0\delta'}\lambda^{(1-\eta)/2-s_*}Z^4\right).
\]
The fourth moment of \(Z\) therefore gives the averaged bound
\[
 O_\diamond(x^{5-C_0\delta'}\lambda^{(1-\eta)/2-s_*}),
\]
which saves a power when
\(C_0\delta'<p_*((1-\eta)/2-s_*)\), with room for the remaining
losses. We choose \(\delta'\) this small, also relative to \(\kappa\),
and then choose the direct Taylor degree.

Next expand the deterministic quotient coefficients from \(K\) to
\(K_S\), keeping the same total truncation degree. These scalar Taylor
remainders have bounded common coefficients and at least \(J_0\) pair
factors in all. Each factor is either \(K_S-K=bf+d\) or \(b\delta\).
Split the pressure-gradient test into its deterministic and linear
stochastic parts as well. We need only the difference between the two
ensembles, so Lemma~\ref{exp:lem-centered-count} applies once its
weight condition is checked. Every designated factor has its envelope
bound, and the density has prefactor \(\lambda\). With \(j\) factors
at \(l\) distinct vertices, changing to the normalized measure costs
\(w^l\). Since \(j+l\ge5\),
\[
 \prod_s dv_s=w^l\prod_s dv_s/w
       \lesssim w^{5-j}\prod_s dv_s/w.
\]
Thus the remainder has the weight \eqref{eq:E1} and is bounded by
\eqref{eq:E2}.

The retained terms recombine by Lemma~\ref{found:reference-change}.
Differentiating a scalar moment toward \(K_S\) inserts
\((K_S-K)\tau_a\tau_b/2\). Allocation symmetry permits these
insertions within products, with each scalar moment evaluated on its
own subtopology. Insertions into a unit moment sum to zero when
attached to the observed topology. Use these identities in the
unrestricted allocation sums, before dividing terms into mass bands
for estimates. The quotient retained through degree \(J_0-1\) is
therefore the quotient at \(K\) with argument
\(J=(K_S-K)+b\delta=b\Delta+d\).

By \eqref{eq:A9}, differentiating the logarithm through degree \(J_0\)
produces this gradient. Any additional highest-degree coefficient
derivatives have \(J_0\) factors and an \(O_\circ(\lambda)\) weight;
after splitting the factors into deterministic and stochastic parts,
the same calculation puts them in \eqref{eq:E1}. Pure deterministic
terms cancel. This establishes \eqref{eq:E3} on the whole remaining
direct path, including masses at which a pointwise small-insertion
argument was unavailable.

For size changes, first consider terms with a base pair at mass at
least \(V_2\). Call a base factor \emph{high} if its mass is at least
\(V_1\); recall that \(V_1<V_2\). Count the square term as two base
factors on the same pair. Expanding \(\Delta=f+\delta\), every high
deterministic factor, including \(d\) or \(\dot d\), is bounded by
\(O_\diamond(x^{2+b_H}(Z^2+Z_g))\). Stochastic high factors have the
variance bounds \eqref{eq:D15}; all factors also have the all-moment
smallness \eqref{eq:D11}. Low factors use their envelopes.

The scalar coefficients are bounded, including \(\dot b=O(1)\).
Differentiating \(K\) may add one scalar pair with weight
\(|\dot K|\). At high mass this weight is integrable separately. If
it is tied to a high base vertex, integrate against that same variable
and use the pointwise bounds \eqref{eq:D11}, \eqref{eq:D15}. At a low
mass \(v\), the inverse estimates give \(|\dot K(v)|=O(v)\).
The terms with at least one stochastic factor fall into the following
cases.
\begin{enumerate}
\item If there are two high stochastic factors, Cauchy--Schwarz gives
\(D_{H,1}^2\). If there is one high stochastic factor and a high
deterministic factor, the bound is
\(D_{H,1}x^{2+b_H}(Z^2+Z_g)\).
\item If every high factor is deterministic, extract the bound for
one of them and apply \eqref{eq:D16} to the low factors. A sole low
factor at \(v\) has the additional scalar coefficient \(O(v)\)
required there: immediately after its split, its two ends lie in
groups with odd parity, as in the proof of
Proposition~\ref{prop:statistics-high}. The square term cannot have
this configuration. An added dotted scalar pair preserves the
parity argument if it splits later than \(v\); if it splits earlier
or at \(v\), its own weight is \(O(v)\).
\item If a stochastic factor is the only high factor, it is the
factor at mass at least \(V_2\), so it costs \(D_{H,2}\). With one low
factor, the preceding parity argument gives
\(\int_{v<V_1}v e(v)\,dv\lesssim_\diamond V_1Y^2\). With two or more
low factors, integrate one using
\(\int e(v)\,dv\lesssim_\diamond Y^2\), and use \eqref{eq:D11} on
another to gain \(x^{\delta_*/2}\). Higher fixed moments and H\"older
give the same bound when these factors share a vertex.
\end{enumerate}
Here is the resulting power count, before \(\gamma\)-averaging and up
to the permitted losses. In the last two rows use
\(Y^2=x^2\lambda^{-s_*}Z^2\) and \(\lambda=x^{p_*}\).
\[
\begin{array}{l|l}
\text{selected factors}&\text{bound}\\ \hline
\text{two stochastic high}&x^{6-2C_Ha_1}Z^4\\
\text{stochastic and deterministic high}
 &x^{5+b_H-C_Ha_1}(Z^4+Z_gZ^2)\\
\text{only deterministic high}
 &x^{5+b_H-Cp_*}(Z^4+Z_gZ^2)\\
\text{sole stochastic high, one low}
 &x^{5+a_1-C_Ha_2-p_*s_*}Z^4\\
\text{sole stochastic high, two or more low}
 &x^{5+\delta_*/2-C_Ha_2-p_*s_*}Z^4.
\end{array}
\]
The hierarchy in \eqref{eq:D15} makes every exponent strictly larger
than five: it gives \(b_H\gg C_Ha_1\), \(a_1\gg C_Ha_2\), and permits
\(p_*\) to be chosen smaller than all remaining gaps. The factors
\(Z^4\) and \(Z_gZ^2\) have bounded \(\gamma\)-expectation, the latter
by Cauchy--Schwarz. Thus every displayed cost has the required
averaged saving. Products with no stochastic factor cancel outright.

On a direct path, terms whose maximum mass is bounded below by a
fixed positive constant already satisfy \eqref{eq:E1}. The dotted
weights supply \(O_\diamond(\lambda)\), and
\(|\dot d|\lesssim_\diamond\lambda r^2e\). For a quadratic term with
two distinct designated vertices, scalar parity gives the required
smaller-scale ratio. If differentiation adds a scalar pair at a
smaller mass \(v'\), its weight
\(|\partial_{\log\lambda}K(v')|\lesssim\lambda v'\) gives that ratio
instead.

All direct terms of degree \(j\ge5\) also satisfy \eqref{eq:E1} by
the normalized-measure calculation above. The same holds for the
remaining size terms of those degrees: now \(w\lesssim V_2\ll\lambda\)
and
\(\prod_s dv_s\lesssim\lambda w^{5-j}\prod_s dv_s/w\).
An extra dotted weight that is unbounded must be at a higher mass and
on a separate variable, where its absolute integral is available.
It remains to treat only low-region terms of degrees \(2,3,4\).

\begin{proposition}[Coefficients after compression]\label{exp:prop-low-coefficients}
Fix the dyadic restrictions on the base factors and a cut \(c\gg w\)
within the region where the small inverse is regular. A sufficiently
large fixed separation factor is enough; on size paths require also
\(c\lesssim\lambda\). Group base slots into blocks that share their
path through \(c\). Their split tree below \(c\) is the \emph{skeleton},
and every designated pair joins distinct blocks.

Apply the operator \({\cal O}\) defined above to the pressure
polynomial \eqref{eq:E3}: it is \(1+\partial_\theta\), with the
formal action \(\partial_\theta\Delta=-\Delta/6\), on a size path,
and \(\partial_{\log\lambda}\), holding \(\Delta\) fixed, on a
direct path. Group the resulting signed allocation coefficients by
this skeleton and sum their internal placements within each block.
The bound below concerns these differentiated, grouped coefficients.
This is a rearrangement of a signed measure. The probabilistic
expectations used to evaluate
the coefficients remain positive scalar laws on each fixed topology.
The spin-model pair tests have been projected to \(c\), or to a
neighboring smaller cut, so they depend only on the skeleton and can
be kept outside the internal sums. For \(j=2,3,4\), the resulting
coefficient, including the skeleton allocation weight, is bounded by
\[
 O_\diamond(\lambda c^{5-j})\prod_{s\ {\rm skeleton\ fork}} dv_s/c .       \tag{E4}\label{eq:E4}
\]
If \(j=2\), the designated vertices are distinct, and \(c\asymp w\),
the bound includes the further factor
\(O_\diamond(\min_s r_s/c)\). A factor \(d\) or \(\dot d\) may first
contribute \(\lambda c^2\) to the coefficient, leaving a deterministic
factor bounded by the ordinary envelope because \(r\lesssim c\).
Like \(b\) and \(\dot b\), these weights are constant during each
internal block sum.

Combine contributions only when they multiply the same product of
pair tests on the same skeleton. In particular, this combines a
repeated pair with the square counterterm. Give such terms the same
dyadic restrictions and projection cuts. The cut, shape, and mass
coordinates are fixed during coefficient differentiation. Internal
summation may then be done before or after that differentiation.
The spin tests receive only the formal action on \(\Delta\) already
specified: project the differentiated polynomial, without
differentiating a random test or a projection law.
\end{proposition}

\paragraph{Reduction to Proposition~\ref{exp:prop-low-coefficients}.}
Start with a sufficiently small fixed cut on direct paths, choosing
the low-mass threshold smaller still, and with a cut of order
\(\lambda\) on size paths. By
Corollary~\ref{cor:statistics-low-projection}, changing one pair to its
projection at \(c\) costs \(O_{p,\circ}(z_c^0c)\). Conditional on the
positive-law prefixes at that cut, a cross-block pair has exactly
its projected conditional mean. Expand a product as projected factors
plus changes. Terms with just one change have zero conditional mean,
so every error contains at least two changes.

For this first projection, bounded coefficient densities suffice.
Their prefactor is \(\lambda\) on direct paths and \(O_\circ(1)\),
including integrated dotted weights, on size paths. The respective
costs are \(O_\diamond(x^6\lambda^{-\eta})\) and
\(O_\diamond(x^6\lambda^{-C})\). Now reduce the cut by factors of two
until it is comparable to \(w\), retaining a large enough separation
factor. For a step from \(c\) to \(c/2\), first sum the coefficients
at \(c\) and apply \eqref{eq:E4}. The product error again contains at
least two changes of size \(O_{p,\circ}(z_c^0c)\); every other factor,
after extracting the defect weights allowed in the proposition, is
\(O_{p,\diamond}(c)\). Inactive masked pairs do not change. Hence the
step costs
\[
 O_\diamond(\lambda c^5(z_c^0)^2)
       =O_\diamond(x^6\lambda^{-\eta}/c)
       \lesssim_\diamond x^5\lambda^{s_*/2-\eta}
\]
because \(c\gtrsim Y\ge U\). At the final cut, \eqref{eq:E4} is
precisely \eqref{eq:E1}: the designated fork measures have the required
normalization, and the remaining skeleton variables have bounded
volumes in units of \(c\). Apply \eqref{eq:E2} to their expectation
difference. All cuts are common to the two ensembles. There are only
fixedly many degrees and shapes, and the dyadic bands and successive
cuts cost subpower factors. Thus Proposition~\ref{exp:prop-low-coefficients}
implies the remaining bounds for \eqref{eq:I}, uniformly also in the
site count, with a strict power saving.


\subsection{Scalar compression at a fixed cut}

\begin{proof}[Proof of Proposition~\ref{exp:prop-low-coefficients}]
We first evaluate the internal signed sum in \eqref{eq:E3} before
applying \({\cal O}\), and then justify commuting that operator with
the sum. Absolute values are taken only after the resulting
cancellations. Fix a cut \(c\) in the regular inverse region with
all designated pair splits strictly below it and
\(c\gtrsim w\ge Y\ge U\); on size paths also \(c\lesssim\lambda\).
Throughout coefficient differentiation, hold this cut, the skeleton
shape, its mass coordinates, and all mass restrictions fixed.

\paragraph{The objects in the compression formula.}
For a degree \(j\) term, partition its \(2j\) labeled endpoint slots
into blocks at \(c\). Let \(D\) be the number of blocks, ordered by
their first slots. The \(j\) ordered pairs form a graph on these
blocks; every edge has two distinct ends. The skeleton allocation
measure is the signed per-root measure for \(D\) ordered distinct end
labels in the hierarchy ending at \(c\), divided by \(c^D\).
Its absolute weight on each of finitely many shapes is bounded by
\[
 C_jc^{-1}\prod_{s\ {\rm skeleton\ fork}}dv_s/c.
\]
Indeed the first label has unit per-root weight, and each of the
remaining \(D-1\) placements contributes scale \(c\). Division by
\(c^D\) leaves \(c^{-1}\). After rescaling mass to \([0,1]\), the
remaining densities are bounded. Attachments at existing forks may
reuse their mass coordinates, but block identities remain distinct.
There are no prescribed old leaves or forks. Thus distinct free fork
coordinates are unequal almost everywhere, while a fork may have
more than two children.

Next let \(\pi\) be a partition of the \(j\) edges into cumulant
parts, keeping repeated edges as separately labeled slots. For a
part \(P\in\pi\), let \(k_{A,P}\) be the number of its endpoints in
block \(A\), and put
\[
 D_\pi=\sum_{P\in\pi}\#\{A:k_{A,P}>0\}.
\]
Thus \(D\) counts blocks in the full skeleton, whereas \(D_\pi\)
counts a block again each time it is used by another part. Neither
counts endpoint multiplicities.

On the fixed skeleton, let \(\zeta^A_u\), \(0\le u\le c\), be the
scalar prefix fields under the positive scalar law. They share their
Brownian driver until their skeleton separation, and thereafter use
independent transitions. Each part \(P\) uses a separate copy of this
law restricted to its blocks. The normalized source sum for \(k\)
slots inside one block is
\[
 s_k(z)=c^{-1}e^{-cV(t(c),z)}\partial_z^k e^{cV(t(c),z)}.
\]
With the explicit edge weights of \eqref{eq:E3} kept outside, the
coefficient multiplying the skeleton measure is
\[
 \sum_{\pi}(-1)^{|\pi|-1}(|\pi|-1)!\ c^{D_\pi}
       \prod_{P\in\pi}\mathbf E_{\rm low}
           \prod_{A\ {\rm in}\ P} s_{k_{A,P}}(\zeta^A_c).
                                                               \tag{E5}\label{eq:E5}
\]
The expectation \(\mathbf E_{\rm low}\) is probabilistic; the preceding
grouping of allocation coefficients by skeleton is a signed sum.

For example, take two copies of the edge \(AB\) between cut blocks.
Their original terminal labels may differ; projection makes their
spin-model pair tests identical. Abbreviate
\(s_1(\zeta^A_c)=\mathsf M_A\),
\(s_2(\zeta^A_c)=\mathsf P_A\). The one-part partition uses two
blocks, whereas the two-part partition uses each block twice. Their
combined internal coefficient is therefore
\[
 c^2\mathbf E_{\rm low}(\mathsf P_A\mathsf P_B)
   -c^4\bigl(\mathbf E_{\rm low}(\mathsf M_A\mathsf M_B)\bigr)^2.
\]
The skeleton measure contributes \(c^{-1}\), giving net powers
\(c\) and \(c^3\). There are two orientations of the second edge
for a fixed first edge, so its factor in \eqref{eq:E3} is
\(2/(2^2 2!)=1/4\). The square counterterm uses this same skeleton,
with internal coefficient \(c^2\) and weight \(-b/4\). These
normalizations will give the quadratic cancellation below.

\paragraph{The finite signed identity.}
First expand each cumulant into products of positive scalar moments
on its indicated edge sets. Once the skeleton assignments are fixed,
there is no restriction on the internal placements within a block.
A part observes only its own induced genealogy, so the internal
signed sums for different parts can be performed separately, even
when their actual terminal leaves coincide.

This factorization can be checked directly at finitely many levels
including \(c\). In the labeled product rule of
\cite[Proposition~2.5]{PE}, the coefficients choosing distinct children
are falling factorials of incoming-to-child mass ratios, and every
terminal source slot carries its terminal mass factor. Keep the
induced genealogies as discrete marks and first take the successive
ratios, including an auxiliary root ratio, to be sufficiently large
positive integers. The coefficients then count child identities.
After distinct block identities have been assigned at \(c\), the
choices for disjoint sets of slots inside a block are independent
counts, with terminal slots chosen with replacement. Their counts
therefore multiply. The resulting rational coefficient identity
extends by polynomial identity to the real masses, keeping terminal
mass fixed and dividing by the root mass before its zero limit.
This argument concerns only the coefficients; it does not extend
spin tests or probability laws from integer ratios.

\paragraph{Evaluation of the internal source sums.}
For one part, the positive scalar prefix law on its restricted
skeleton is independent of the internal placements after \(c\).
Conditional on these prefix fields, the source product rule sums
\(k\) internal slots in a block to
\(e^{-cV}\partial_z^k e^{cV}=c s_k\). Different blocks use independent
future kernels. Thus every used block contributes one factor \(c\),
giving \(c^{D_\pi}\), and the cumulant partition formula gives
\eqref{eq:E5}. Signed path limits and bounded fixed-order derivatives,
as supplied by PE, pass this finite-level identity to the continuous
paths. Equivalently, one may fix the prefix fields and approximate
only the future clocks within the blocks. There are no prescribed
internal split times to preserve in that approximation.

\paragraph{Differentiation at a fixed skeleton.}
The cut and skeleton measure stay fixed. Differentiating the scalar
moments adds one covariance pair with weight \(\dot K\). By
\eqref{eq:P8}, its absolute integrated weight is finite; its
non-atomic placements have bounded density for fixed cut and label
count. The same bound holds along affine comparison chords of the
clocks, so the finite signed differentiation rule and Fubini apply
to the internal sums.

To pass from chords to ordinary derivatives, fix \(N\). The
clock difference quotients are uniformly bounded in mass: for size
paths this follows from the smoothing floor and clipping, and for
direct paths from the localized smooth inverse perturbation. For
almost every parameter they converge at every mass in the low
region and at almost every other mass. The scalar kernels converge
by fixed-size path continuity. Dominated convergence therefore
passes the divided covariance insertions to the dotted insertions,
also for almost every fixed skeleton. A reused split is either one
of these integration coordinates or a prescribed low skeleton mass,
where convergence is pointwise; constant diagonals still cancel.
This proves that internal summation and a.e. coefficient
differentiation commute. We may choose the low-region partition
separately at each parameter when estimating the result: the
coefficient has already been differentiated, so no moving indicator
or cut is being differentiated.

\paragraph{Scalar jets and common-prefix estimates.}
Write
\(\chi=V_{zz}(0,0)\), \(g=V_{zzzz}(0,0)\),
\(k_0=t'(0+)>0\), and \(G(u,z)=V(t(u),z)\).
The notation \(O_*(A)\) controls both the value and the parameter
derivative: on size paths the derivative is \(\partial_\theta\),
and on direct paths it is \(\lambda^{-1}\partial_{\log\lambda}\).
The same convention applies to fixed \(L^p\) norms when fields are
coupled by common Brownian drivers in mass coordinates. Derivatives
of the explicit weights \(b,d\) will be handled separately.

By \eqref{eq:P1}--\eqref{eq:P2}, the inverse \(t\) is smooth here,
with \(t'\) bounded above and below and with the required higher
derivatives bounded. The positive-order spatial derivatives of \(G\)
also have bounded mass and spatial derivatives of every required
fixed order, allowing polynomial growth in \(z\) after parameter
differentiation. To see the mass derivatives, use
\[
 \partial_uG=-\tfrac12t'(G_{zz}+uG_z^2).
\]
For direct derivatives, differentiate the scalar PDE in \(\lambda\)
at fixed normalized time, as in the lift proof, and then compose with
the regular inverse. For size derivatives, use the rescaled scalar
jets, including the terminal-mass derivative already estimated in
Section~\ref{sec:paths}. These operations preserve evenness in the
field. In particular, on size paths,
\[
 \partial_\theta k_0=-3 k_0/6,\qquad
 \partial_\theta\chi=\chi/6+O(e^{-\Omega(M)}),\qquad
 \partial_\theta g=3g/6+O(e^{-\Omega(M)}).                              \tag{E6}\label{eq:E6}
\]
The normalized inverse slope and initial spatial derivatives are
constant under this rescaling, up to the exponentially small terminal
effect. This gives \eqref{eq:E6}. Moreover, \(\chi\) is uniformly
bounded above and away from zero. At every designated split,
\[
 b=b_\circ+O(\lambda c^2),\qquad b_\circ=\chi^{-2},
\]
and the error has the same bound under unscaled parameter
differentiation. By \eqref{eq:D2}, each \(d\) or \(\dot d\) can
likewise contribute \(O_\diamond(\lambda c^2)\) to the coefficient,
leaving its ordinary envelope factor.

The low prefix equation is
\[
 d\zeta=\sqrt{t'}\,dB+u t'G_z(u,\zeta)du.
\]
Starting from zero, this equation gives
\(\zeta_u=O_*(\sqrt u)\) and
\(\zeta_u-\sqrt{k_0}B_u=O_*(u^{3/2})\) in every fixed moment.
For a conditional evolution starting at \(z\), the field and its
parameter derivative have norms at most
\(C_p(|z|+\sqrt c)(1+|z|)^C\). Its spatial flow derivatives have
bounded moments, allowing polynomial growth in \(z\). These bounds
follow by differentiating the additive-noise equation: the drift is
smooth with bounded spatial derivatives, the differentiated
coefficients have polynomial growth, and each odd differentiated
drift vanishes at zero.

Evaluate the block observables at their fields at \(c\), unless
another time is specified. Taylor expansion, together with even/odd
parity and the definition of \(s_k\), gives
\[
\begin{split}
 \mathsf M_A := s_1(\zeta^A_c)&=G_z(c,\zeta^A_c)=\chi\zeta^A_c+O_*(c^{3/2}),\\
 \mathsf P_A:=s_2(\zeta^A_c)&=G_{zz}(c,\zeta^A_c)+c(G_z(c,\zeta^A_c))^2
                        =\chi+O_*(c),\\
 s_3(\zeta^A_c)&=g\zeta^A_c+O_*(c^{3/2}),\qquad
 s_4(\zeta^A_c)=g+O_*(c).
\end{split}                                                            \tag{E7}\label{eq:E7}
\]
For blocks \(A,B\) splitting at \(v_{AB}<c\), the martingale
property of \(G_z\) reduces their product expectation to the common
prefix at \(v_{AB}\). Thus
\[
 \mathbf E_{\rm low} \mathsf M_A\mathsf M_B=\chi^2 k_0 v_{AB}+O_*(v_{AB}^2).
\]
We also need
\[
 \mathbf E_{\rm low} \mathsf P_A\mathsf P_B
   =\chi^2+\chi^3 k_0 c^2+\tfrac12 g^2 k_0^2 v_{AB}^2+O_*(c^3).       \tag{E8}\label{eq:E8}
\]
To obtain \eqref{eq:E8}, keep the terminal cut \(c\) fixed and set
\(H_u=G_{zz}(u,\zeta_u)+cG_z(u,\zeta_u)^2\). The PDE and It\^o's
formula give drift \((c-u)t'G_{zz}^2\) and diffusion coefficient
\(\sqrt{t'}(G_{zzz}+2cG_zG_{zz})\). In the product of the two block
processes, the two drift integrals sum to
\(\chi^3k_0c^2+O_*(c^3)\), by \eqref{eq:E7}. Their common quadratic
variation stops at \(v_{AB}\); its expected rate before that time is
\(g^2k_0^2u+O_*(cu)\). Its integral is the term
\(\tfrac12g^2k_0^2v_{AB}^2\), with an error of order \(O_*(c^3)\).

For three blocks we use
\[
 \mathbf E_{\rm low} \mathsf P_A \mathsf M_B\mathsf M_C=\chi^3 k_0 v_{BC}+O_*(c^2).
                                                                        \tag{E9}\label{eq:E9}
\]
When the designated splits are distinct, this estimate retains the
earlier mass. Suppose \(v_{AB}=h'<v_{AC}\), so that \(v_{BC}=h'\).
Condition at \(h'\) on the common field \(z\). After subtracting
\(\chi\mathsf M_B\mathsf M_C\), the conditional expectation on the
\(A,C\) side is that of \((\mathsf P_A-\chi)\mathsf M_C\). It is an
odd function of \(z\), with spatial derivative bounded by
\(C(c+z^2)(1+|z|)^C\), also after parameter differentiation at fixed
\(z\). Indeed \(\mathsf P-\chi=O(c+|\zeta_c|^2)\), its spatial
derivative is \(O(|\zeta_c|)\), and \(\mathsf M=O(|\zeta_c|)\)
with bounded derivative; apply the conditional flow bounds and
H\"older to their product.

The \(B\)-side conditional mean is \(G_z(h',z)\), of size
\(O_*(|z|(1+|z|)^C)\). Substituting the common prefix field
\(z=O_*(\sqrt{h'})\) therefore bounds the error in \eqref{eq:E9}
by \(O_*(ch')\). The other ordering is identical after exchanging
\(B,C\). Without a designated earlier split, \eqref{eq:E7} gives
the stated \(O_*(c^2)\) error.

\paragraph{The coefficients canceled by size rescaling.}
For a pure degree \(j\) term with no defect factors, the formal
rule \(\partial_\theta\Delta=-\Delta/6\) makes the size operator on
its coefficient equal to \(1-j/6+\partial_\theta\). Since
\(\ell=e^{\theta/6}\),
\[
 \left(1-\frac j6+\partial_\theta\right)
       \ell^{-(6-j)}=0.
\]
Equation~\eqref{eq:E6} and \(b_\circ=\chi^{-2}\) identify exactly
the coefficients with this scaling:
\[
\begin{array}{c|c|c}
 j&\text{coefficient monomials}&\text{size scaling}\\ \hline
 2&b_\circ^2\chi^3k_0,\quad b_\circ^2g^2k_0^2&\ell^{-4}\\
 3&b_\circ^3g^2k_0,\quad b_\circ^3\chi^3&\ell^{-3}\\
 4&b_\circ^4g^2&\ell^{-2}.
\end{array}
\]
The masses multiplying these monomials are held fixed. Exponentially
small terminal errors will be retained until the end. On direct
paths, the operator is only \(\partial_{\log\lambda}\), and the
\(O_*\) convention already supplies the factor \(\lambda\).
For a term with a defect, use the product rule and extract at least
one factor \(\lambda c^2\) as above. In the remaining case analysis
we include the net factor \(c^{D_\pi-1}\); the normalized skeleton
fork measure in \eqref{eq:E4} is left outside.

\Needspace{8\baselineskip}
\paragraph{The low-degree cases.}
\begin{enumerate}
\item Suppose \(j=2\) and \(D=2\). The two edges join the same
blocks, so use the repeated-edge calculation following \eqref{eq:E5}.
Its coefficient combines with the square counterterm because both
multiply the same projected pair square. After extracting the common
skeleton measure and \(c^2/4\), their coefficient before
differentiation is
\[
 b^2\left\{\mathbf E_{\rm low}\mathsf P_A\mathsf P_B
       -c^2(\mathbf E_{\rm low}\mathsf M_A\mathsf M_B)^2\right\}-b.
\]
The constant term \(b^2\chi^2-b\), and its unscaled derivatives, are
\(O(\lambda c^2)\) because \(b_\circ\chi^2=1\) and
\(b-b_\circ=O(\lambda c^2)\). The next two terms from \eqref{eq:E8},
after replacing \(b\) by \(b_\circ\), are
\[
 b_\circ^2\chi^3k_0c^2,
 \qquad \tfrac12b_\circ^2g^2k_0^2v_{AB}^2.
\]
Their direct derivatives are \(O(\lambda c^2)\); their size
contributions cancel by the \(j=2\) row of the table. Substitution
errors from \(b-b_\circ\) have the required bound. The remaining
\(O_*(c^3)\) errors, including the smaller product-of-means term,
contribute \(O(c^4)\) after the net compression factor \(c\).
On size paths this is \(O(\lambda c^3)\) because \(c\lesssim\lambda\);
on direct paths the derivative supplies \(\lambda\) itself.
Thus \eqref{eq:E4} holds. A defect factor supplies
\(\lambda c^2\) directly and gives the same bound, with the allowed
\(O_\diamond(1)\) loss.

\item Suppose \(j=2\) and \(D=3\), with edges \(AB,AC\).
For the one-part partition, \eqref{eq:E9} gives the leading
coefficient \(b_\circ^2\chi^3k_0v_{BC}\), with net compression
factor \(c^2\). Its size contribution cancels by the same row of
the table; its direct derivative costs \(O(\lambda c^3)\), or
\(O(\lambda c^2h')\) when \(v_{BC}=h'\) is the earlier split.
The error is \(O_*(c^2)\), or \(O_*(ch')\) when the designated
vertices are distinct and their earlier mass is \(h'\).
Multiplication by \(c^2\), followed on size paths by
\(c\lesssim\lambda\), gives respectively
\(O(\lambda c^3)\) and \(O(\lambda c^2h')\). The errors from
\(b-b_\circ\) and any defect factors satisfy these same bounds.
For the two-part partition the product of pair moments is
\(O_*(v_{AB}v_{AC})\), with net factor \(c^3\), which is smaller
still. Since \(h'\lesssim\min_s r_s\), the refined bound gives the
extra ratio required in \eqref{eq:E4}.

\item Suppose \(j=2\) and \(D=4\). Both partitions have net
factor \(c^3\), and their products of moments are \(O_*(c^2)\).
When the designated vertices are distinct with earlier mass \(h'\),
these products are also \(O_*(h')\). For a product of two pair
moments this follows at once from their individual bounds. For the
four-spin moment, condition on the prefixes at \(h'\). The earlier
edge has separated ends, while the ends of the later edge are still
together. Hence at least two of the independent future groups have
odd parity. Each such group's conditional expectation vanishes at
zero starting field and has a polynomially bounded spatial
derivative, also after parameter differentiation. The common prefix
fields have size \(O_*(\sqrt{h'})\), so the product of these two
odd expectations gives \(O_*(h')\).

On size paths the refined contribution is therefore
\(O(c^3h')\le O(\lambda c^2h')\), exactly the required smaller-scale
bound. On direct paths the derivative supplies \(\lambda\).
The unrefined \(O_*(c^2)\) estimate gives \eqref{eq:E4} without the
ratio. Extracting defect factors preserves both conclusions.

\item Suppose \(j=3\), for which the target is
\(O_\diamond(\lambda c^2)\). Terms with defects or with
\(D_\pi\ge4\) already have enough powers of \(c\): on a size
path one extra power is bounded by \(\lambda\), and on a direct
path the derivative supplies \(\lambda\). It remains to consider
one-part partitions with two or three blocks.
For \(D=2\), each block has degree three and \eqref{eq:E7} gives
\(g^2k_0v_{AB}+O_*(c^2)\). For \(D=3\), an odd block degree gives
\(O_*(c)\); otherwise all three degrees are two, giving
\(\chi^3+O_*(c)\). The two leading monomials after multiplication
by \(b_\circ^3\) are precisely those in the \(j=3\) row of the
table, so their size contributions cancel. With net factors \(c\)
and \(c^2\), respectively, the remaining errors are \(O(c^3)\),
which is at most \(O(\lambda c^2)\). Direct derivatives give the
same target without this cancellation.

\item Suppose \(j=4\), for which the target is
\(O_\diamond(\lambda c)\). Only size terms with no defect and
\(D_\pi=2\) lack the required power before cancellation. Such a
term has one part, on two blocks of degree four each. Its moment is
\(g^2+O_*(c)\). Multiplication by \(b_\circ^4\) gives the last
row of the table, so the leading size contribution cancels. The
remainder with net factor \(c\) is \(O(c^2)\), at most
\(O(\lambda c)\). All other terms have an additional compression
power, a defect factor, or the direct derivative factor
\(\lambda\).
\end{enumerate}
The exponentially small terminal errors retain the mass factors of
their corresponding main terms. Since \(M=\log^3(1/x)\), they are
smaller than every fixed power of \(x\). This proves \eqref{eq:E4}.
\end{proof}

The row expansion, the high-mass estimates, and the low-degree
coefficient bound now give both increments required in
Proposition~\ref{prop:increment}. The final parameter choice in
Section~\ref{sec:parameters} makes their savings simultaneous and
justifies the stated uniformity for moving site weights.
